Este capítulo apresenta como a solução foi testada e avaliada. A Seção~\ref{sec:planejamento-testes} descreve o planejamento. A Seção~\ref{sec:testes-automatizados} relata os testes automatizados do núcleo de tradução e do serviço de \textit{back-end} e sua execução na integração contínua. A Seção~\ref{sec:avaliacao-especialistas} trata do teste de usabilidade com o painel de especialistas, e a Seção~\ref{sec:avaliacao-estudantes}, da comparação entre duas versões de um programa feita com uma turma de estudantes. Em cada seção, o resultado de cada teste é discutido logo depois de apresentado. As Seções~\ref{sec:resultado-literatura}, \ref{subsec:limites-interpretacao} e~\ref{sec:resultado-sintese} relacionam os resultados aos trabalhos anteriores, delimitam o que eles permitem concluir e os confrontam com os objetivos.

\section{Planejamento dos testes}
\label{sec:planejamento-testes}

O planejamento seguiu a pirâmide de testes apresentada na Subseção~\ref{subsec:testes-tdd}: muitos testes pequenos e rápidos na base e menos testes nos níveis que exercitam o sistema inteiro. Cada nível ficou com um objeto e uma ferramenta:

\begin{itemize}
    \item Testes do núcleo de tradução: verificam os \textit{tokens}, cada \textit{scanner} do analisador léxico e a aceitação ou rejeição de programas inteiros pela gramática, sob diferentes configurações da linguagem. Usam o Vitest e estão na Subseção~\ref{subsec:testes-nucleo}.
    \item Testes do serviço de \textit{back-end}: verificam as rotas do FastAPI, os contratos de resposta e as regras de acesso, com um banco de dados de teste. Usam o pytest e estão na Subseção~\ref{subsec:testes-backend}. Foram escritos com o auxílio do Claude Code, conforme a Seção~\ref{sec:uso-ia}.
    \item Integração contínua: as suítes rodam no GitHub Actions, e a implantação só ocorre se forem aprovadas (Subseção~\ref{subsec:ci-backend}).
    \item Teste de usabilidade com especialistas: verifica a percepção de professores e profissionais de computação sobre o núcleo de tradução e a plataforma, por meio de um formulário com escalas consolidadas, o SUS, o UEQ-S e o TAM, e de blocos escritos para o projeto. Está na Seção~\ref{sec:avaliacao-especialistas}.
    \item Comparação A/B com estudantes: verifica qual de duas versões do mesmo programa, uma na sintaxe padrão e outra personalizada em português, os alunos de uma turma declaram entender com mais facilidade. Está na Seção~\ref{sec:avaliacao-estudantes}.
\end{itemize}

\section{Testes automatizados}
\label{sec:metodologia-testes}\label{sec:testes-automatizados}

Os testes automatizados verificam o núcleo de tradução e o serviço de \textit{back-end} e rodam também na integração contínua.

\subsection{Testes do núcleo de tradução}
\label{subsec:testes-nucleo}

A verificação automatizada das fases do compilador utiliza, na versão do repositório descrita neste texto, o \textit{framework} Vitest. Sua adoção pode ser conferida no arquivo \texttt{packages/compiler/package.json}, que declara a dependência e o comando \texttt{test}, e nas importações de \texttt{describe}, \texttt{it} e \texttt{expect} dos arquivos de teste. A suíte está organizada no diretório \texttt{packages/}\allowbreak\texttt{compiler}\allowbreak\texttt{/src/tests} e estruturada de modo a refletir a topologia interna do projeto. Os testes foram subdivididos em três camadas: testes de \textit{tokens}, que verificam a estabilidade dos identificadores numéricos e dos agrupamentos por categoria; testes de \textit{lexer}, que exercitam cada \textit{scanner} especializado em isolamento, abrangendo comentários, literais numéricos, cadeias com escapes, terminadores configuráveis, delimitadores de bloco, operadores em forma de palavra e geração de \textit{tokens} de indentação; e testes de gramática, que validam a aceitação ou rejeição de programas inteiros sob distintas configurações de tipagem, modos de bloco e variantes do comando \texttt{switch}.

A adoção do TDD discutida por \citeonline{aniche2012} mostrou-se particularmente útil em construções como o terminador de instrução configurável e o modo de blocos por indentação, em que pequenas alterações na lógica do \textit{lexer} poderiam quebrar inadvertidamente programas que antes compilavam corretamente. O ciclo iterativo do TDD, composto pelas etapas de escrita do teste, implementação mínima e refatoração, conforme ilustra a Figura \ref{fig:ciclo_tdd}, é a referência metodológica. Em conformidade com a pirâmide de testes apresentada na Seção~\ref{subsec:testes-tdd}, a maior parte da suíte concentra-se no nível unitário, garantindo execuções rápidas e \textit{feedback} imediato, enquanto os testes de gramática operam como testes de integração internos ao compilador, verificando o comportamento agregado das fases léxica, sintática e de geração de código intermediário.

Para tornar concreta essa estratégia, o Trecho de Código~\ref{lst:teste-aliases} reproduz um dos casos que verificam a aplicação de um mapeamento configurável de operadores no analisador léxico. O teste instancia o \texttt{Lexer} com um mapeamento que substitui os operadores lógicos por palavras e verifica que os \textit{tokens} produzidos carregam exatamente os mesmos identificadores de tipo que seriam gerados pelos símbolos canônicos. Essa equivalência é o que garante que as fases seguintes do compilador permaneçam indiferentes à personalização escolhida pelo usuário.

\begin{lstlisting}[language=typescript, caption={Caso de teste da equivalência entre operadores simbólicos e operadores em forma de palavra (\texttt{operator-word-aliases.spec.ts}).}, label={lst:teste-aliases}, captionpos=b]
describe("operator word aliases", () => {
  it("tokenizes logical aliases with the same token IDs as symbols",
     () => {
    const lexer = new Lexer("if a and not b {}", {
      operatorWordMap: {
        logical_and: "and",
        logical_not: "not",
      },
    });

    const tokens = lexer.scanTokens();

    expect(tokens.map((token) => token.type)).toEqual(
      expect.arrayContaining([
        TOKENS.LOGICALS.logical_and,
        TOKENS.LOGICALS.logical_not,
      ]),
    );
  });
});
\end{lstlisting}


Na verificação local de 19 de outubro de 2026, a suíte do núcleo aprovou 243 testes em 22 arquivos (Apêndice~\ref{ap:repositorio}). Em 19 de setembro eram 224 testes em 21 arquivos (Apêndice~\ref{ap:repositorio}). Os casos acrescentados cobrem incremento e decremento, declarações múltiplas com inicializador no fim da lista, leitura de variáveis não declaradas, o limite de passos, delimitadores informados como símbolos, literais numéricos no fim do código-fonte e o resumo de regras em \texttt{language-rules.spec.ts}. Os casos incluem também igualdade comum e estrita, ternários aninhados com execução somente do ramo selecionado, avisos de conversão e rejeição de atribuições e retornos incompatíveis. O Apêndice~\ref{ap:testes} apresenta a distribuição e o comando utilizado.

\subsection{Testes do serviço FastAPI}
\label{subsec:testes-backend}

A suíte do serviço Python utiliza pytest e pytest-asyncio, com modo assíncrono automático definido na configuração do projeto. Os cenários combinam verificações de esquemas com testes das rotas e de sua integração com serviços e persistência. As \textit{fixtures} centralizam a preparação da aplicação e do banco, e as fábricas de dados criam usuários, organizações e atividades necessários a cada cenário. Essa organização corresponde aos mecanismos de preparação reutilizável apresentados na Seção~\ref{subsec:fundamentos-testes-api} \cite{pytest2026fixtures,pytestasyncio2026} (Apêndice~\ref{ap:repositorio}).

O ambiente emprega SQLite em memória e sessões assíncronas do SQLAlchemy. Uma fixture cria as tabelas a partir dos modelos e insere a organização e o usuário de sistema; ao final da suíte, remove as tabelas. Os testes que utilizam o banco recebem uma sessão com transação e reversão ao término. A dependência de sessão da aplicação é substituída por essa sessão, e o HTTPX utiliza transporte ASGI para enviar requisições às rotas no próprio processo. Assim, as verificações observam respostas HTTP e efeitos no banco de teste, sem depender do serviço PostgreSQL da implantação \cite{fastapi2026async,fastapi2026overrides} (Apêndice~\ref{ap:repositorio}).

Entre os casos relacionados à segurança estão o cadastro com retorno dos dados do usuário sem senha, a recusa de credenciais inválidas, a restrição de acesso à submissão de outro aluno e a impossibilidade de um aluno atribuir notas. Também há verificações de ausência de propriedades nas respostas resumidas: os testes de apresentação de linguagens conferem a omissão da configuração completa, enquanto os testes do histórico do exercício verificam que códigos enviados e submissões de colegas não sejam expostos naquele retorno (Apêndice~\ref{ap:repositorio}).

A Figura~\ref{lst:resumo-submissao} reproduz parte de um caso existente. A preparação da turma, dos alunos e das submissões foi omitida da listagem e está disponível no arquivo original. A asserção do estado HTTP precede a inspeção do conteúdo, de modo que uma resposta de erro sem os campos verificados não seja confundida com um resultado correto (Apêndice~\ref{ap:repositorio}).

\begin{figure}[htbp]
\centering
\caption{Recorte do teste que verifica a ausência de código no histórico resumido do exercício.}
\label{lst:resumo-submissao}
\begin{lstlisting}[language=Python]
response = await async_client.get(
    f"/exercises/{exercise.id}?listId={exercise_list.id}",
    headers={"Authorization": f"Bearer {token}"},
)

assert response.status_code == 200, response.text
for submission in response.json()["submissions"]:
    assert "codeSnapshot" not in submission
assert "alice-v1" not in response.text
\end{lstlisting}
\legend{Fonte: elaborado pelo autor.}
\end{figure}

\subsection{Resultados e alcance da verificação}
\label{subsec:resultados-testes-backend}

Em 19 de setembro de 2026, a execução local da suíte completa do back-end aprovou 188 casos em 17 arquivos, sem falhas ou casos ignorados, em aproximadamente 59 segundos. Foram utilizados Python 3.12.12, pytest 9.0.2, pytest-asyncio 1.3.0, HTTPX 0.28.1 e FastAPI 0.135.1. A Tabela~\ref{tab:testes-backend} agrupa os arquivos por finalidade, contando separadamente os casos parametrizados coletados pelo executor. O registro estruturado da revisão apresenta a distribuição por arquivo e os comandos utilizados (Apêndice~\ref{ap:repositorio}).

\begin{table}[htbp]
\centering
\caption{Distribuição dos testes executados no back-end.}
\label{tab:testes-backend}
\small
\renewcommand{\arraystretch}{1.15}
\begin{tabular}{|>{\raggedright\arraybackslash}p{0.28\textwidth}|c|c|>{\raggedright\arraybackslash}p{0.43\textwidth}|}
\hline
\textbf{Grupo} & \textbf{Arquivos} & \textbf{Casos} & \textbf{Exemplos de verificação} \\\hline
Autenticação, usuários e organizações & 3 & 29 & Cadastro, login, identificação e operações de perfil. \\\hline
Turmas, exercícios e listas & 3 & 44 & Criação, edição, participação e organização de atividades. \\\hline
Políticas de linguagem e registros de configuração & 2 & 39 & Precedência da linguagem obrigatória e configuração preservada na submissão. \\\hline
Linguagens & 4 & 44 & Operações do acervo, esquemas, apresentação e carga inicial. \\\hline
Submissões e histórico próprio & 2 & 15 & Envio, reenvio, correção, acesso e campos do resumo. \\\hline
Paginação & 1 & 4 & Páginas, totais, busca e compatibilidade das listagens. \\\hline
Cenários complementares & 2 & 13 & Saúde do serviço, erros de autenticação e operações recusadas. \\\hline
\textbf{Total} & \textbf{17} & \textbf{188} & Casos coletados e aprovados na revisão local. \\\hline
\end{tabular}
\legend{Fonte: elaboração própria a partir do registro de verificação de 19 de setembro de 2026 (Apêndice~\ref{ap:repositorio}).}
\end{table}
\FloatBarrier

A contagem corresponde a casos executados e não a uma porcentagem de cobertura do código ou das políticas de segurança. Como as tabelas são criadas diretamente pelos modelos e a dependência de sessão é substituída, essa execução não verifica as migrações Alembic, o comportamento específico do PostgreSQL nem o encerramento da transação da dependência de produção. A aprovação dos cenários existentes também não elimina as lacunas identificadas na Seção~\ref{subsec:limites-seguranca-backend} (Apêndice~\ref{ap:repositorio}).

\subsection{Testes no fluxo de integração contínua}
\label{subsec:ci-backend}

O workflow atual executa verificações nas solicitações de integração e nos envios à branch \texttt{main}. O job denominado \texttt{unit} reúne Vitest para compilador e IDE e pytest para o back-end; seu nome não implica que todos os casos sejam unitários, pois há integração de rotas e persistência. Outro job executa Playwright com serviços locais em Docker. A implantação por SSH depende da aprovação dos dois jobs e só ocorre em envio à \texttt{main}. Esse arquivo documenta a condição para deploy (Apêndice~\ref{ap:repositorio}).

Nos merges examinados, a dependência de sessão passou a encerrar a transação antes do envio da resposta, utilizando o escopo de função do FastAPI. A resposta de autenticação também passou a incluir o modelo de usuário, evitando uma consulta imediata adicional para obter seus dados. Essas alterações relacionam a consistência do serviço ao fluxo observado pela interface; sua presença no código é distinta da evidência produzida pela suíte com SQLite. A documentação do FastAPI descreve o momento de encerramento das dependências com esse escopo \cite{fastapi2026yield} (Apêndice~\ref{ap:repositorio}).


\section{Teste de usabilidade com especialistas}
\label{sec:protocolo-avaliacao}\label{sec:avaliacao-especialistas}

O teste de usabilidade foi feito com um formulário respondido por um painel de especialistas, chamado aqui de formulário dos especialistas. O formulário é compartilhado com o texto do coautor: foi enviado ao mesmo grupo de respondentes, e cada trabalho analisa os itens que interessam ao seu tema. Aqui interessam os itens sobre o núcleo de tradução. O formulário recebeu oito respostas entre 1º e 9 de novembro de 2026 e está reproduzido no Apêndice~\ref{ap:formulario-especialistas}. As primeiras subseções descrevem o instrumento e o procedimento; as seguintes apresentam os resultados de cada bloco, já com a discussão. O resumo agregado das respostas está no Apêndice~\ref{ap:respostas}.

\subsection{Quem responde e o que é possível concluir}

O formulário dos especialistas foi respondido por um painel formado por professores e profissionais de computação. Neste trabalho, considera-se especialista quem atua ou já atuou com programação, como docente de disciplinas da área ou como profissional de desenvolvimento de software. Não é um estudo com os alunos, que são o público-alvo da plataforma. Pediu-se ao respondente que avaliasse se a ferramenta é adequada para um aluno em primeiro contato com programação, e não que relatasse a própria experiência de aprendizagem.

Essa escolha limita o que o trabalho pode afirmar. Com um painel de especialistas é possível dizer que profissionais da área acham plausível que a personalização da sintaxe reduza a barreira enfrentada pelo iniciante. Não é possível dizer que ela reduz. Verificar isso exige medir o desempenho de estudantes, o que está registrado na Seção~\ref{subsec:limites-interpretacao} como etapa seguinte. O formulário dos estudantes não muda esse limite, porque também colhe opinião e não mede desempenho.

Para o tema deste trabalho, porém, o painel dá o que um aluno iniciante não teria como dar: uma avaliação técnica sobre se o subconjunto implementado da linguagem é suficiente, se a representação intermediária exibida corresponde à execução real e se as mensagens de erro continuam claras depois que o usuário troca o vocabulário da linguagem.

\subsection{Composição do instrumento}

O formulário foi respondido pelo próprio participante e distribuído pela internet. Ele tem doze blocos, sete deles com grades de itens em escala, como mostra a Tabela~\ref{tab:blocos-pa01}. Três blocos vêm de instrumentos já consolidados na literatura de avaliação de interfaces. O primeiro é a \textit{System Usability Scale} (SUS), de \citeonline{brooke1996}, que mede a facilidade de uso percebida com dez afirmações e resume as respostas em um escore de 0 a 100. O segundo é a versão curta do \textit{User Experience Questionnaire} (UEQ-S), de \citeonline{schrepp2017}, que avalia a impressão geral sobre o produto por meio de pares de adjetivos opostos. O terceiro vem do \textit{Technology Acceptance Model} (TAM), de \citeonline{davis1989}, cujas afirmações tratam da utilidade percebida e da intenção de uso. Os outros nove blocos foram escritos para este projeto, porque nenhum instrumento pronto pergunta sobre o que se investiga aqui. Os formatos de resposta e o cálculo de cada escala estão na Subseção~\ref{subsec:escalas-escores}.

\begin{table}[htbp]
\centering
\caption{Blocos do formulário dos especialistas, procedência e número de itens.}
\label{tab:blocos-pa01}
\footnotesize
\renewcommand{\arraystretch}{1.2}
\begin{tabular}{|p{5.2cm}|p{4.6cm}|c|}
\hline
\textbf{Bloco} & \textbf{Procedência} & \textbf{Itens} \\\hline
Termo de consentimento (texto de leitura, sem item de marcação) & Elaborado para o projeto & -- \\\hline
Perfil profissional & Elaborado para o projeto & 4 \\\hline
Uso em sala (condicional) & Elaborado para o projeto & 10 \\\hline
Como a plataforma foi usada & Elaborado para o projeto & 4 \\\hline
Avaliação dos recursos & Elaborado para o projeto & 10 \\\hline
Facilidade de uso & SUS \cite{brooke1996} & 10 \\\hline
Impressão geral & UEQ-S \cite{schrepp2017} & 8 \\\hline
Potencial pedagógico & Elaborado para o projeto & 8 \\\hline
Riscos e limitações & Elaborado para o projeto & 3 \\\hline
Qualidade técnica & Elaborado para o projeto & 6 \\\hline
Utilidade e intenção de uso & TAM \cite{davis1989} & 6 \\\hline
Comentários abertos & Elaborado para o projeto & 5 \\\hline
\end{tabular}
\legend{Fonte: elaborado pelo autor.}
\end{table}

A versão aplicada difere em dois pontos da que foi descrita na revisão anterior deste texto. A afirmação sobre a suficiência do subconjunto da linguagem saiu do bloco de qualidade técnica, e a afirmação sobre alunos da mesma turma usarem sintaxes diferentes saiu do bloco de riscos. As duas foram para o bloco de uso em sala, que só aparece para quem declara lecionar ou já ter lecionado. Esse bloco ficou com nove afirmações em escala e uma pergunta sobre o momento da disciplina em que a plataforma seria mais útil. A consequência é que os dois itens deslocados foram respondidos por no máximo quatro pessoas, e não pelo painel inteiro.

O SUS, o UEQ-S e o TAM foram usados porque têm valores de referência publicados. Com eles é possível comparar o escore obtido com o de outros produtos avaliados, em vez de apenas descrever a opinião dos respondentes. Antes do bloco de avaliação dos recursos, o formulário mostra o mesmo programa em duas versões: uma na configuração padrão, com \texttt{int}, \texttt{for}, \texttt{if} e delimitadores simbólicos, e outra personalizada, com \texttt{numero\_inteiro}, \texttt{para}, \texttt{se}, operadores escritos por extenso e os delimitadores \texttt{inicio} e \texttt{fim}. O par mostra, em um caso concreto, a equivalência entre vocabulários diferentes que o núcleo de tradução garante.

\subsection{Itens sobre o núcleo de tradução}

O bloco de qualidade técnica é o que trata mais de perto do tema deste trabalho. Ele é dirigido a quem tem formação na área e traz a opção ``não sei''. As afirmações são:

\begin{itemize}
    \item as mensagens de erro são claras e indicam o que precisa ser corrigido;
    \item a ferramenta respondeu de forma rápida e estável durante o uso;
    \item o assistente de criação de linguagem é compreensível para quem não tem formação em compiladores;
    \item a validação prévia impede adequadamente configurações de linguagem inválidas;
    \item o código intermediário exibido corresponde ao que o programa realmente executa.
\end{itemize}


Esse último item merece atenção. Ele coloca em avaliação externa uma propriedade que o texto sustenta pela forma como o sistema foi construído: a representação intermediária mostrada ao usuário é a mesma que o interpretador consome, e não uma ilustração feita à parte. Uma avaliação negativa nesse item indicaria diferença entre o que a interface exibe e o que a máquina executa, o que comprometeria a premissa central do trabalho.

O bloco termina com uma questão de múltipla escolha sobre construções que faltam na linguagem, com alternativas para registros, orientação a objetos, tratamento de exceções, manipulação de cadeias mais completa, bibliotecas gráficas e divisão do programa em vários arquivos, além da alternativa que diz não ter faltado nada relevante para o ensino introdutório. Essa questão delimita, com base nas respostas, até onde vai o subconjunto implementado e ajuda a priorizar os trabalhos futuros no núcleo. A afirmação de que o subconjunto, com tipos, condicionais, repetição, arranjos, funções e entrada e saída, é suficiente para uma disciplina introdutória ficou no bloco de uso em sala, com a opção ``não se aplica''.

Do bloco de avaliação dos recursos, que traz a opção ``não usei'', interessam três itens, todos sobre a exibição das fases intermediárias da tradução: se a tabela de \textit{tokens} ajuda a entender como a plataforma lê o código, se o código intermediário ajuda a entender o que o programa faz e se a execução passo a passo ajuda a acompanhá-lo. Soma-se o item sobre as mensagens de erro ajudarem a identificar o que corrigir, que mede o resultado prático de adaptar essas mensagens ao vocabulário em uso. O bloco de potencial pedagógico repete os três primeiros sob outro ângulo, perguntando pelo valor didático deles para quem está aprendendo, e acrescenta que mostrar os erros no vocabulário escolhido pelo próprio aluno facilita a correção. Essa repetição é proposital: separa o que o especialista achou útil para si, usando a ferramenta, do que ele considera útil para um iniciante.

\subsection{Escalas usadas e cálculo dos escores}
\label{subsec:escalas-escores}

Os blocos em escala usam dois formatos de resposta. O SUS e os blocos de uso em sala, avaliação dos recursos, potencial pedagógico, riscos, qualidade técnica e utilidade e intenção de uso pedem o grau de concordância com cada afirmação em cinco pontos, em que 1 corresponde a ``discordo totalmente'', 3 ao ponto neutro e 5 a ``concordo totalmente''. Nos blocos elaborados para o projeto há também uma opção fora da escala, que não entra nas médias: ``não se aplica'' no uso em sala, ``não usei'' na avaliação dos recursos e ``não sei'' nos blocos de potencial pedagógico, de riscos e de qualidade técnica. O UEQ-S usa sete pontos entre dois adjetivos opostos. As perguntas de perfil, as de uso da plataforma, a do momento da disciplina e a das construções que fizeram falta são de múltipla escolha. O último bloco tem quatro questões abertas e um item de múltipla escolha em que o respondente diz se autoriza a citação das suas respostas sem identificação.

A facilidade de uso usa a escala de \citeonline{brooke1996}, com dez afirmações que alternam entre positivas e negativas. Essa alternância faz parte do instrumento e não deve ser corrigida. Para calcular o escore, subtrai-se 1 do valor marcado nos itens ímpares, subtrai-se de 5 o valor marcado nos itens pares, somam-se os resultados e multiplica-se o total por 2,5. O resultado vai de 0 a 100 e não é uma porcentagem. Nos estudos reunidos por \citeonline{sauro2012}, a média dos escores fica perto de 68 pontos. Para interpretar o escore, \citeonline{bangor2009} propõem faixas de aceitabilidade: abaixo de 50 pontos o resultado é considerado inaceitável, de 50 a 70 fica na faixa marginal e acima de 70 é considerado aceitável. A impressão geral usa a versão curta do questionário de \citeonline{schrepp2017}, com oito pares de adjetivos opostos em sete pontos; subtrai-se 4 do valor bruto e separam-se a dimensão pragmática e a hedônica. Como no instrumento original, o termo negativo de cada par fica sempre à esquerda e o positivo à direita, e os quatro itens pragmáticos vêm antes dos quatro hedônicos \cite{schrepp2023handbook}. A orientação fixa é própria da versão curta: na versão completa, de 26 itens, metade dos pares começa pelo termo positivo. A utilidade e a intenção de uso vêm do modelo de \citeonline{davis1989}.

O bloco de riscos funciona ao contrário dos outros: nele, concordar significa apontar um problema. As afirmações são:

\begin{itemize}
    \item aprender em uma sintaxe personalizada pode dificultar a migração posterior para uma linguagem de mercado;
    \item a liberdade de personalizar pode desviar a atenção do aluno do raciocínio algorítmico;
    \item o tempo gasto configurando a linguagem não compensa o ganho obtido.
\end{itemize}
 A afirmação sobre alunos da mesma turma usando sintaxes diferentes, deslocada para o bloco de uso em sala, tem o mesmo sentido e é lida junto com essas três. Os escores desses itens são relatados separadamente e não podem ser somados aos do bloco de potencial pedagógico, porque as duas medidas apontam em sentidos opostos e uma anularia a outra. O bloco foi incluído de propósito: se o instrumento trouxesse apenas afirmações favoráveis, a aprovação obtida diria muito pouco.

\subsection{Participantes, procedimento e questões éticas}

Os participantes do painel foram convidados diretamente e não são identificados. O formulário registra titulação, área de atuação, ferramentas de ensino que a pessoa já conhece e tempo de experiência com programação. Em uma avaliação feita por especialistas, esses dados fazem parte do resultado e são apresentados junto com os escores. Um item pede que o respondente marque o que fez na plataforma e o que chegou a personalizar, o que permite separar a avaliação de quem usou a ferramenta da de quem não experimentou o recurso avaliado.

Sobre o número de participantes: no modelo de \citeonline{nielsen1993}, cada novo participante revela uma parte menor de problemas ainda não encontrados, e por isso um grupo pequeno, de cinco a oito pessoas, já revela a maior parte dos problemas de usabilidade e basta para analisar as respostas abertas. Para apresentar o escore SUS com intervalo de confiança são necessários vinte respondentes ou mais. Responderam oito pessoas. Por isso o escore é apresentado como indicativo, sem intervalo de confiança e sem afirmar que existe diferença entre subgrupos.

Os dois formulários foram aplicados como pesquisa de opinião com participantes não identificados e não passaram pelo comitê de ética da instituição. O formulário dos especialistas abre com um termo de consentimento, apresentado como texto de leitura, sem item de marcação, e por isso as respostas desse formulário não trazem registro individual de consentimento. No formulário dos estudantes, o termo de assentimento é um item de resposta obrigatória, e as dezoito respostas trazem a opção ``Li e quero participar''. A sessão com estudantes envolveu menores de idade: nove tinham até 15 anos e nove tinham 16 ou 17. O formulário registra o assentimento de cada aluno, mas não a autorização de um responsável. Essas lacunas no registro do consentimento entram como limitação do estudo no Capítulo~\ref{cap:consideracoes}.

\subsection{Leitura dos resultados}
\label{subsec:leitura-resultados}

Os números apresentados a seguir vêm do resumo das respostas, gerado por um \textit{script} a partir das planilhas de respostas (Apêndice~\ref{ap:respostas}). Os blocos de análise, a ordem e os critérios de leitura foram definidos antes da coleta e foram mantidos; defini-los depois de ver os dados permitiria escolher o recorte mais favorável à solução. Os critérios de cálculo dos escores estão na Subseção~\ref{subsec:escalas-escores}.

As convenções de leitura valem para todas as subseções de resultados. Com oito respondentes, o texto usa contagens, médias e medianas, sem porcentagens. Nas afirmações em escala de 1 a 5, as respostas ``não usei'', ``não sei'' e ``não se aplica'' ficam fora da média e são contadas à parte. Nenhuma resposta foi excluída. Quando é preciso apontar um respondente, ele é indicado por R1 a R8, na ordem de chegada das respostas.

\subsection{Composição do painel}
\label{sec:resultado-painel}

Esta subseção usa as questões de perfil (\ref{qE:perfil-ini} a \ref{qE:perfil-fim}) e de uso da plataforma (\ref{qE:uso-ini} a \ref{qE:uso-fim}) do Apêndice~\ref{ap:formulario-especialistas}. Os oito respondentes foram convidados diretamente e formam uma amostra de conveniência. Não são identificados. Pela maior titulação declarada, três têm doutorado (a alternativa do formulário era ``doutorado ou superior''), dois têm graduação e três têm ensino médio, sem curso superior concluído. Cinco atuam como profissionais de desenvolvimento de software e três são professores do ensino superior, dois dos quais também lecionam no ensino técnico. Quatro lecionam ou já lecionaram disciplina que envolva programação: os três professores e um dos profissionais de desenvolvimento. Dois trabalham com programação há menos de 2 anos, três de 2 a 5 anos, dois de 6 a 10 anos e um há mais de 10 anos.

Sobre as ferramentas de ensino já conhecidas, cinco citaram o Portugol Studio ou o Portugol Webstudio, quatro citaram Scratch, Blockly ou similar, três citaram Beecrowd, URI ou outro juiz \textit{online}, dois citaram o VisuAlg, dois citaram Thonny, Python Tutor ou outro visualizador de execução, e um não conhecia nenhuma delas.

O tempo de uso declarado foi curto para metade do painel: quatro usaram a plataforma por menos de 15 minutos, um de 15 a 30 minutos, um de 30 a 60 minutos e dois por mais de 1 hora. A Tabela~\ref{tab:resultado-atividades} mostra o que os respondentes disseram ter feito e o que personalizaram.

\begin{table}[htbp]
\centering
\caption{Atividades declaradas pelos respondentes do painel e elementos da linguagem que personalizaram (respostas múltiplas, oito respondentes).}
\label{tab:resultado-atividades}
\footnotesize
\renewcommand{\arraystretch}{1.2}
\begin{tabular}{|p{10.5cm}|c|}
\hline
\textbf{O que fez na plataforma} & \textbf{Respondentes} \\\hline
Criou ou editou uma linguagem no assistente & 5 \\\hline
Escreveu e executou um programa & 5 \\\hline
Viu o código intermediário & 4 \\\hline
Usou a execução passo a passo & 4 \\\hline
Resolveu e submeteu um exercício & 3 \\\hline
Viu a tabela de \textit{tokens} & 2 \\\hline
Entrou em uma turma ou criou uma turma & 2 \\\hline
Importou uma linguagem da comunidade & 1 \\\hline
\textbf{O que personalizou na linguagem} & \textbf{Respondentes} \\\hline
Palavras dos tipos & 6 \\\hline
Modo de tipagem & 5 \\\hline
Entrada e saída & 4 \\\hline
Operadores em palavra & 3 \\\hline
Controle de fluxo & 2 \\\hline
Literais booleanos & 2 \\\hline
Delimitadores de bloco ou indentação & 2 \\\hline
Terminador de instrução & 1 \\\hline
Usou uma linguagem predefinida, sem alterar & 2 \\\hline
Não personalizou nada & 2 \\\hline
\end{tabular}
\legend{Fonte: dados da pesquisa (2026).}
\end{table}

Para este trabalho importa quem chegou às fases intermediárias da tradução. Quatro respondentes declararam ter visto o código intermediário (R1, R2, R4 e R5), quatro declararam ter usado a execução passo a passo (R1, R3, R4 e R5) e só dois declararam ter visto a tabela de \textit{tokens} (R4 e R5). A tabela de \textit{tokens} foi, portanto, a fase menos visitada. Do lado da personalização, as palavras dos tipos e o modo de tipagem foram os elementos mais alterados, e o terminador de instrução foi alterado por um respondente.

Três características das respostas condicionam a leitura das subseções seguintes. A primeira é o pouco contato de parte do painel. R7 não marcou nenhuma atividade nem elemento personalizado, declarou menos de 15 minutos de uso e deu as notas mais baixas em vários blocos. R8 marcou apenas os recursos de turma e de exercício e escreveu, nas questões abertas, que não conseguiu explorar a plataforma. Os dois lecionam. Isso não desqualifica a opinião deles, mas indica que parte das notas reflete uma primeira impressão.

A segunda é o padrão de R2, que marcou o mesmo valor em todos os itens de cada bloco em escala: 3 no bloco de uso em sala, no bloco de recursos e no SUS, 6 nos oito pares do UEQ-S e 4 nos blocos de potencial pedagógico, de riscos, de qualidade técnica e de utilidade. As respostas de R2 foram mantidas, e o efeito delas é indicado onde pesa.

A terceira é a relação entre o que foi avaliado e o que foi usado. Em vários itens o número de notas é maior que o número de pessoas que declararam ter usado o recurso. Um dos dois respondentes que marcaram não ter personalizado nada marcou também três elementos personalizados, e a resposta foi mantida como veio.

\subsection{Avaliação do núcleo de tradução}
\label{sec:resultado-nucleo}

Esta subseção trata das questões de qualidade técnica (\ref{qE:qual-ini} a \ref{qE:qual-fim}) e das questões de avaliação dos recursos sobre as fases da tradução (\ref{qE:fases-ini} a \ref{qE:fases-fim}) do Apêndice~\ref{ap:formulario-especialistas}. A Tabela~\ref{tab:resultado-nucleo} reúne os cinco itens em escala do bloco de qualidade técnica e os três itens do bloco de avaliação dos recursos que tratam da exibição das fases intermediárias. No primeiro bloco a resposta à parte é ``não sei''; no segundo, ``não usei''.

\begin{table}[htbp]
\centering
\caption{Itens sobre o núcleo de tradução (escala de 1 a 5; ``não sei'' ou ``não usei'' contados à parte).}
\label{tab:resultado-nucleo}
\footnotesize
\renewcommand{\arraystretch}{1.2}
\begin{tabular}{|p{7.8cm}|c|c|c|c|p{1.3cm}|}
\hline
\textbf{Afirmação} & \textbf{n} & \textbf{Média} & \textbf{Mediana} & \textbf{Mín.--máx.} & \textbf{À parte} \\\hline
\multicolumn{6}{|l|}{\textbf{Qualidade técnica (``não sei'' à parte)}} \\\hline
As mensagens de erro são claras e indicam o que precisa ser corrigido & 7 & 4,71 & 5 & 4--5 & \centering\arraybackslash 1 \\\hline
A ferramenta respondeu de forma rápida e estável durante o uso & 8 & 4,88 & 5 & 4--5 & \centering\arraybackslash 0 \\\hline
O assistente de criação de linguagem é compreensível para quem não tem formação em compiladores & 8 & 4,38 & 4,5 & 3--5 & \centering\arraybackslash 0 \\\hline
A validação prévia impede adequadamente configurações de linguagem inválidas & 6 & 4,67 & 5 & 4--5 & \centering\arraybackslash 2 \\\hline
O código intermediário exibido corresponde ao que o programa realmente executa & 6 & 4,83 & 5 & 4--5 & \centering\arraybackslash 2 \\\hline
\multicolumn{6}{|l|}{\textbf{Avaliação dos recursos (``não usei'' à parte)}} \\\hline
A tabela de \textit{tokens} é útil para entender como a plataforma lê o código & 6 & 4,50 & 5 & 3--5 & \centering\arraybackslash 2 \\\hline
O código intermediário é útil para entender o que o programa faz & 7 & 4,14 & 4 & 3--5 & \centering\arraybackslash 1 \\\hline
A execução passo a passo é útil para acompanhar o programa & 7 & 4,57 & 5 & 3--5 & \centering\arraybackslash 1 \\\hline
\end{tabular}
\legend{Fonte: dados da pesquisa (2026).}
\end{table}

Nenhum item do bloco de qualidade técnica recebeu nota abaixo de 3, e quatro dos cinco não receberam nota abaixo de 4. O item mais bem avaliado foi a resposta rápida e estável da ferramenta, com média 4,88: sete respondentes marcaram 5 e R2 marcou 4. Todos tinham condição de julgar esse item, porque ele trata do uso em geral, e não de um recurso específico. A validação prévia da configuração teve média 4,67 entre seis respondentes, e o assistente teve média 4,38, com uma nota 3, dada por R8, que não abriu o assistente.

O item sobre o código intermediário exibido corresponder ao que o programa executa recebeu destaque no protocolo, porque uma avaliação negativa nele indicaria diferença entre o que a interface mostra e o que o interpretador executa. Não houve avaliação negativa. Seis respondentes deram nota, com média 4,83: cinco marcaram 5 e R2 marcou 4. Dois marcaram ``não sei'', R7 e R8, que não declararam ter visto o código intermediário. Num bloco dirigido a quem tem formação na área, essa resposta indica que o recurso não foi visto.

O número de notas, porém, é maior que o de pessoas que declararam ter visto o recurso. Dos seis que avaliaram o item, quatro declararam ter visto o código intermediário (R1, R2, R4 e R5). R3 e R6 deram nota 5 sem ter marcado essa atividade. Assim, o julgamento sobre a correspondência entre o código exibido e a execução se apoia, de forma declarada, em quatro respondentes, um deles com o padrão de R2. Mesmo entre esses quatro, o formulário não permite saber se a pessoa comparou o código intermediário com a execução ou se avaliou a impressão geral do recurso. O item mede percepção; a correspondência em si decorre de como o sistema foi construído: a lista de instruções exibida na IDE é a mesma que o interpretador percorre, como descrito na Subseção~\ref{subsec:geracao-codigo-intermediario}, e é verificada pelos testes da Seção~\ref{sec:metodologia-testes}.

Os três itens sobre as fases intermediárias ficaram acima de 4. A execução passo a passo teve média 4,57, com cinco notas 5. A tabela de \textit{tokens} teve média 4,50 entre seis respondentes, com dois ``não usei'', de R6 e R8. O código intermediário teve a menor média dos três, 4,14, com três notas 5, dois 4 e dois 3; as notas 3 vieram de R2 e de R7. Também aqui o número de notas supera o de usos declarados: a tabela de \textit{tokens} recebeu seis notas, mas só dois respondentes declararam tê-la visto.

O bloco de potencial pedagógico faz a mesma pergunta sob outro ângulo, sobre o valor didático de cada fase para quem está aprendendo. A execução passo a passo teve média 4,75, o código intermediário 4,43 e a tabela de \textit{tokens} 4,29, os dois últimos entre sete respondentes, com um ``não sei'' de R8. As médias para o próprio uso e para o iniciante ficaram próximas, sem uma direção comum: o código intermediário subiu de 4,14 para 4,43, e a tabela de \textit{tokens} desceu de 4,50 para 4,29. A diferença mais clara está em R7, que deu nota 4 à tabela de \textit{tokens} e 3 ao código intermediário como recursos e nota 2 aos dois como valor didático para o iniciante. Esse respondente julgou as duas fases úteis para si e pouco úteis para um iniciante.

\subsection{Suficiência do subconjunto da linguagem}
\label{sec:resultado-subconjunto}

Esta subseção trata da afirmação sobre a suficiência do subconjunto (questão \ref{qE:subconjunto}) e da pergunta sobre as construções que fizeram falta (questão \ref{qE:qual-fim}) do Apêndice~\ref{ap:formulario-especialistas}. A afirmação de que o subconjunto da linguagem, com tipos, condicionais, repetição, arranjos, funções e entrada e saída, é suficiente para uma disciplina introdutória ficou no bloco de uso em sala, como registrado na Seção~\ref{sec:protocolo-avaliacao}. Dos quatro que lecionam, três deram nota e um marcou ``não se aplica''. As notas foram 5, 5 e 3, com média 4,33. O 3 é de R2. Com três notas, o resultado é um indício.

A pergunta sobre construções que fizeram falta foi feita a todo o painel e teve sete respondentes; R4 a deixou em branco. Cinco marcaram não ter sentido falta de nada relevante para o ensino introdutório. Registros ou estruturas compostas foram marcados por dois respondentes, e modularização em vários arquivos, orientação a objetos e manipulação de cadeias mais rica por um cada. Tratamento de exceções e bibliotecas gráficas não foram marcados. Um respondente, R7, marcou ao mesmo tempo que não sentiu falta de nada relevante e que sentiu falta de registros e de orientação a objetos, e a resposta foi mantida como veio. Se for lida só pelas construções marcadas, quatro respondentes não apontaram falta alguma.

O painel não convergiu para uma construção ausente. A mais citada, registros ou estruturas compostas, teve duas marcações, e nenhuma outra passou de uma. Orientação a objetos e modularização em vários arquivos já constam das limitações do Capítulo~\ref{cap:consideracoes}. Registros e manipulação de cadeias mais rica não constavam e passam a constar. A linguagem tem o tipo \texttt{string} e literais de cadeia, mas o formulário não detalha que operações o respondente esperava. Para priorizar os trabalhos futuros no núcleo, o resultado indica registros como o candidato mais citado, sem que duas marcações bastem para tratá-lo como demanda do painel.

\subsection{Mensagens de erro com vocabulário personalizado}
\label{sec:resultado-diagnosticos}

Esta subseção trata das questões \ref{qE:qual-ini}, \ref{qE:msg-recursos}, \ref{qE:msg-ped} e \ref{qE:msg-aberta} do Apêndice~\ref{ap:formulario-especialistas}, todas sobre as mensagens de erro. Três itens em escala tratam das mensagens de erro, um em cada bloco. No bloco de qualidade técnica, a afirmação de que as mensagens são claras e indicam o que precisa ser corrigido teve média 4,71 entre sete respondentes: cinco marcaram 5, R2 e R7 marcaram 4, e R8 marcou ``não sei''. No bloco de avaliação dos recursos, a afirmação de que as mensagens ajudaram a identificar o que corrigir teve média 4,25 entre os oito: quatro marcaram 5, dois marcaram 4 e dois marcaram 3, R2 e R4. No bloco de potencial pedagógico, a afirmação de que relatar os erros no vocabulário escolhido pelo próprio aluno facilita a correção teve média 4,25: cinco marcaram 5, R2 marcou 4, R8 marcou 3 e R7 marcou 2.

As três médias ficaram acima de 4. A questão aberta sobre mensagens confusas traz outro dado. Três respondentes escreveram nela. R7 respondeu que nenhuma foi confusa e R8 que não explorou essa parte. R3 escreveu que ``seria interessante apresentar mensagens de erro mais descritivas, indicando com clareza onde está o problema e, sempre que possível, sugerindo como corrigi-lo''. R3 deu nota 5 aos três itens em escala sobre mensagens de erro. O mesmo respondente, portanto, considera as mensagens claras e úteis e pede que elas sejam mais descritivas e sugiram a correção.

Os itens em escala perguntam se as mensagens indicam o que corrigir, e a resposta aberta pede algo que elas não fazem: propor a correção. As mensagens do núcleo registram um código, o texto, a linha e a coluna do problema, como descrito na Subseção~\ref{subsec:interpretacao}, e nenhuma delas sugere uma correção. O pedido de R3 descreve, portanto, um recurso que o núcleo não tem.

O protocolo previa rastrear cada relato de mensagem confusa até a fase que a emitiu. O relato de R3 não cita uma mensagem específica, e por isso não pode ser rastreado. Nenhum respondente relatou uma mensagem que tenha deixado de acompanhar o vocabulário em uso. Para registrar como as mensagens se comportam nesse ponto, o autor compilou, na versão atual do núcleo, um programa com palavras personalizadas em que falta o ponto e vírgula antes de \texttt{enquanto}, a palavra escolhida para \texttt{while}. A mensagem emitida pela análise sintática foi: ``Token inesperado na linha 1, coluna 32. Esperado `final de instrução', mas recebeu `palavra reservada while' (lexema `enquanto')''. O lexema aparece no vocabulário do usuário, mas a classe do \textit{token} é nomeada pela palavra canônica. Essa mensagem vem do método \texttt{consume} de \texttt{TokenIterator}, que monta o rótulo a partir do identificador do \textit{token}, e não a partir da configuração em uso. A adaptação das mensagens ao vocabulário personalizado é, portanto, parcial. Essa verificação é do autor e não fez parte do instrumento.

\subsection{Usabilidade (SUS) e experiência de uso (UEQ-S)}
\label{sec:resultado-usabilidade}

Esta subseção trata do SUS (questões \ref{qE:sus-ini} a \ref{qE:sus-fim}) e do UEQ-S (questões \ref{qE:ueq-ini} a \ref{qE:ueq-fim}) do Apêndice~\ref{ap:formulario-especialistas}. O escore SUS médio do painel foi 75,6, com mediana 76,25, desvio padrão 16,5, mínimo 50,0 e máximo 95,0. A média fica acima da média de referência de 68 pontos \cite{sauro2012} e acima de 70, valor a partir do qual \citeonline{bangor2009} tratam o resultado como aceitável. Cinco respondentes ficaram acima de 68 e três ficaram abaixo. Com oito respondentes, o resultado é indicativo: não se calcula intervalo de confiança e não se comparam subgrupos, conforme o critério fixado no protocolo. A ressalva sobre a redação em português dos itens, registrada na Seção~\ref{subsec:limites-interpretacao}, também vale para essa comparação.

O escore 50,0 é o de R2, que marcou 3 nos dez itens. Sem essa resposta a média seria 79,3. O valor é informado para que o leitor veja o peso desse padrão; a média relatada continua sendo a dos oito. Entre os itens positivos do SUS, o de menor concordância foi ``acho que gostaria de usar esta plataforma com frequência'', com média 3,50, e o de maior foi ``considerei que as várias funções da plataforma estavam bem integradas'', com 4,25. Os cinco itens negativos tiveram médias entre 1,62 e 1,88.

No UEQ-S, a dimensão pragmática teve média 1,78, a hedônica 1,88 e o valor geral 1,83, na faixa de $-3$ a $+3$. Na média do painel as duas dimensões quase coincidem. R7 teve o único valor negativo, $-0,50$ na dimensão pragmática, com nota 2, em 7, nos pares atrapalhou ou apoiou e ineficiente ou eficiente. Com oito respondentes, essas médias servem só como indicação.

Para o núcleo de tradução, o dado de usabilidade mais direto é o item de qualidade técnica sobre a resposta rápida e estável, com média 4,88, e o par ineficiente ou eficiente do UEQ-S, com média bruta 5,88 em 7. A única nota baixa nesse par foi o 2 de R7.

\subsection{Potencial pedagógico e intenção de uso}
\label{sec:resultado-pedagogico}

Esta subseção trata do bloco de potencial pedagógico (questões \ref{qE:ped-ini} a \ref{qE:ped-fim}) e do bloco de utilidade e intenção de uso, adaptado do TAM (questões \ref{qE:tam-ini} a \ref{qE:tam-fim}), do Apêndice~\ref{ap:formulario-especialistas}. A Tabela~\ref{tab:resultado-pedagogico} reúne o bloco de potencial pedagógico e o bloco de utilidade e intenção de uso. Os dois foram respondidos pelos oito participantes. O que eles registram é a opinião de professores e profissionais sobre o que tende a acontecer com um iniciante, e não uma medida feita com iniciantes.

\begin{table}[htbp]
\centering
\caption{Blocos de potencial pedagógico e de utilidade e intenção de uso (escala de 1 a 5; ``não sei'' contado à parte).}
\label{tab:resultado-pedagogico}
\footnotesize
\renewcommand{\arraystretch}{1.2}
\begin{tabular}{|p{8.0cm}|c|c|c|c|c|}
\hline
\textbf{Afirmação} & \textbf{n} & \textbf{Média} & \textbf{Mediana} & \textbf{Mín.--máx.} & \textbf{Não sei} \\\hline
\multicolumn{6}{|l|}{\textbf{Potencial pedagógico}} \\\hline
Permitir que o aluno defina as palavras-chave tende a reduzir a barreira sintática no primeiro contato & 8 & 3,88 & 4 & 1--5 & 0 \\\hline
Usar palavras em português no lugar de termos em inglês facilita a compreensão do que o código faz & 8 & 4,38 & 5 & 2--5 & 0 \\\hline
A personalização tende a reduzir a frequência de erros de sintaxe cometidos por iniciantes & 7 & 4,14 & 5 & 1--5 & 1 \\\hline
Relatar os erros no vocabulário escolhido pelo próprio aluno facilita a correção do código & 8 & 4,25 & 5 & 2--5 & 0 \\\hline
Expor a tabela de \textit{tokens} tem valor didático para quem está aprendendo & 7 & 4,29 & 5 & 2--5 & 1 \\\hline
Expor o código intermediário tem valor didático para quem está aprendendo & 7 & 4,43 & 5 & 2--5 & 1 \\\hline
A execução passo a passo tem valor didático para quem está aprendendo & 8 & 4,75 & 5 & 4--5 & 0 \\\hline
O conjunto de elementos configuráveis é adequado ao propósito da ferramenta & 7 & 4,57 & 5 & 4--5 & 1 \\\hline
\multicolumn{6}{|l|}{\textbf{Utilidade e intenção de uso}} \\\hline
Esta plataforma é útil para quem está começando a programar & 8 & 4,38 & 5 & 2--5 & 0 \\\hline
Esta plataforma ajudaria alunos iniciantes a aprender lógica de programação mais rápido & 8 & 3,88 & 4,5 & 1--5 & 0 \\\hline
Comparada às ferramentas que eu conheço, esta plataforma oferece algo que elas não oferecem & 8 & 4,62 & 5 & 4--5 & 0 \\\hline
Recomendaria esta plataforma a um colega da área & 8 & 4,25 & 4,5 & 3--5 & 0 \\\hline
Vejo a plataforma como complemento, e não como substituto, à linguagem oficial da disciplina & 8 & 4,38 & 5 & 2--5 & 0 \\\hline
Indicaria esta plataforma a alguém que está começando a programar & 8 & 3,88 & 4,5 & 1--5 & 0 \\\hline
\end{tabular}
\legend{Fonte: dados da pesquisa (2026).}
\end{table}

O protocolo pedia que se discutisse se a opinião do painel concorda com a premissa de concentrar a flexibilidade da linguagem na análise léxica. O formulário não pergunta sobre essa decisão de arquitetura, que fica invisível para quem usa a plataforma. Ele pergunta pelos efeitos que ela deveria produzir. A afirmação de que o conjunto de elementos configuráveis é adequado ao propósito da ferramenta teve média 4,57, com notas 4 e 5 entre os sete que responderam. Esse conjunto é exatamente o que a análise léxica consegue absorver: palavras, operadores em forma de palavra, literais e delimitadores, mais as variantes predefinidas descritas na Seção~\ref{sec:especificacao-linguagem}. Nenhuma resposta aberta pediu outros elementos configuráveis.

Os efeitos esperados tiveram concordância menor. Na afirmação sobre a redução da barreira sintática, seis respondentes concordaram, um ficou no ponto neutro e um discordou totalmente; a média foi 3,88, a menor do bloco. Sobre o português facilitar a compreensão do código, sete concordaram e um discordou, com média 4,38. Sobre a redução dos erros de sintaxe, seis concordaram, um discordou totalmente e um marcou ``não sei'', com média 4,14 entre sete. A maioria do painel concorda com as três premissas, mas não por unanimidade: R7 discordou das três, com notas 1, 2 e 1.

No bloco de utilidade, a afirmação mais aceita foi a de que a plataforma oferece algo que as ferramentas conhecidas pelo respondente não oferecem, com média 4,62 e notas 4 e 5. As duas de menor média, ambas com 3,88, são as que mais se aproximam de uma promessa de resultado: que a plataforma ajudaria iniciantes a aprender lógica mais rápido e que o respondente a indicaria a alguém que está começando. O painel vê a plataforma como complemento à linguagem oficial da disciplina, e não como substituto dela, com média 4,38. Entre os quatro que lecionam, a afirmação ``eu usaria esta plataforma em uma disciplina introdutória de programação'' teve média 3,00, com notas 2, 3, 3 e 4. Os mesmos quatro deram média 4,50 à utilidade de travar uma linguagem em um exercício e à correção automática pelos casos de teste. Os recursos foram bem avaliados por eles, e a adoção não. Na pergunta sobre o momento da disciplina em que a plataforma seria mais útil, dois escolheram ``em paralelo à linguagem oficial, como apoio'', um escolheu ``antes da linguagem oficial da disciplina, como preparação'' e um escolheu ``não vejo espaço para ela na disciplina''. Essa última resposta é de R7, que teve pouco contato com a plataforma. A objeção dele trata do lugar de uma linguagem personalizável na disciplina, questão que o núcleo não resolve, e fica registrada como resultado.

\subsection{Riscos e limitações apontados pelo painel}
\label{sec:resultado-riscos}

Esta subseção trata do bloco de riscos (questões \ref{qE:riscos-ini} a \ref{qE:riscos-fim} do Apêndice~\ref{ap:formulario-especialistas}) e da afirmação de risco do bloco de uso em sala (questão \ref{qE:sala-risco}). Nos itens desta seção, concordar significa apontar um problema, e não aprovar. Eles são relatados separadamente e não entram em nenhuma média junto com o bloco de potencial pedagógico. A Tabela~\ref{tab:resultado-riscos} traz as três afirmações do bloco de riscos e a afirmação de risco do bloco de uso em sala, respondida só por quem leciona.

\begin{table}[htbp]
\centering
\caption{Itens de risco (escala de 1 a 5, em que concordar aponta um problema).}
\label{tab:resultado-riscos}
\footnotesize
\renewcommand{\arraystretch}{1.2}
\begin{tabular}{|p{5.8cm}|c|c|c|c|c|c|}
\hline
\textbf{Afirmação} & \textbf{n} & \textbf{Média} & \textbf{Mediana} & \textbf{Discordam} & \textbf{Neutros} & \textbf{Concordam} \\\hline
Aprender em uma sintaxe personalizada pode dificultar a migração posterior para uma linguagem de mercado & 8 & 3,00 & 3,5 & 3 & 1 & 4 \\\hline
A liberdade de personalizar pode desviar a atenção do aluno do raciocínio algorítmico & 7 & 2,57 & 2 & 4 & 0 & 3 \\\hline
O tempo gasto configurando a linguagem não compensa o ganho obtido & 7 & 1,86 & 2 & 6 & 0 & 1 \\\hline
Alunos de uma mesma turma usando sintaxes diferentes dificultaria o trabalho do professor (só quem leciona) & 3 & 4,00 & 4 & 0 & 1 & 2 \\\hline
\end{tabular}
\legend{Fonte: dados da pesquisa (2026).}
\end{table}

O risco de migração dividiu o painel: quatro concordaram que aprender em uma sintaxe personalizada pode dificultar a passagem para uma linguagem de mercado, um ficou no ponto neutro e três discordaram. Do lado do núcleo, cada palavra personalizada é convertida no mesmo \textit{token} que a palavra canônica, como mostra o Trecho de Código~\ref{lst:teste-aliases}. Por isso todo programa escrito em um vocabulário personalizado tem um equivalente exato no vocabulário padrão. A plataforma, porém, não converte o programa de um vocabulário para outro, e não há nenhum dado sobre a migração, de modo que o risco continua em aberto.

O risco de a personalização desviar a atenção do raciocínio algorítmico teve média 2,57, abaixo do ponto médio da escala, mas três dos sete que responderam concordaram, um deles totalmente. O risco de o tempo de configuração não compensar foi o menos aceito: seis discordaram e só R2 concordou, com média 1,86. Esses dois riscos tratam do uso da personalização em sala, que o núcleo não controla. O que o núcleo oferece é que a linguagem pode vir pronta: as linguagens predefinidas e a configuração travada por exercício chegam ao núcleo como uma configuração a mais, aplicada a cada análise.

A afirmação sobre alunos da mesma turma com sintaxes diferentes foi respondida por três professores, com notas 3, 4 e 5, e nenhum discordou. Essa é a objeção a que a implementação responde de forma mais direta. A rota de validação recebe o código junto com a configuração de linguagem do aluno e executa os casos de teste com ela, como descrito na Seção~\ref{sec:integracao-plataforma}, e a submissão guarda essa configuração no campo \texttt{language\_snapshot}. A correção automática, portanto, não depende de todos usarem a mesma sintaxe. Quando o exercício tem linguagem obrigatória, a interface ainda bloqueia a troca de linguagem. O que a implementação não resolve é a leitura humana: o professor que abre submissões em modo livre continua diante de várias sintaxes.

\subsection{Respostas abertas}
\label{sec:resultado-abertas}

Esta subseção trata das questões abertas \ref{qE:abertas-ini} a \ref{qE:abertas-fim} do Apêndice~\ref{ap:formulario-especialistas}. O formulário tem quatro questões abertas e um item de autorização de citação. Sete respondentes autorizaram a citação sem identificação e um não autorizou. Cinco escreveram em pelo menos uma das questões, todos entre os que autorizaram. As respostas são poucas e curtas, e são relatadas aqui uma a uma.

Sobre o que funcionou melhor, houve quatro respostas. Duas apontam o assistente: ``configuração da linguagem'' e ``criação da linguagem''. Uma aponta os recursos de turma: ``criação de lista e acesso à lista criada''. A quarta, de R3, fala da ``facilidade de utilização'' e da ``possibilidade de visualizar e compreender o funcionamento dos comandos de maneira dinâmica''. Essa última é a única resposta aberta que se aproxima das fases intermediárias expostas pelo núcleo, embora não cite nenhuma delas pelo nome.

Sobre os pontos de dificuldade, também houve quatro respostas. Uma trata da linguagem: segundo R3, ``pode haver alguma dificuldade para compreender a sintaxe e a estrutura dos comandos, especialmente para quem já está habituado a outras linguagens de programação''. É o risco de migração no sentido inverso, de quem já conhece uma linguagem de mercado e encontra outra sintaxe. As outras três tratam da interface ou do uso: a dificuldade de visualizar o botão de criar linguagem, um ``não houve'' e o relato de R8 de que não conseguiu explorar significativamente a plataforma.

A questão sobre mensagens de erro foi discutida na Seção~\ref{sec:resultado-diagnosticos}. Sobre o que mudariam, houve três respostas: um botão de criar linguagem ``mais óbvio/chamativo'', a ampliação dos recursos interativos, ``com mais exemplos práticos, exercícios progressivos e desafios'', e a declaração de R8 de que não explorou o suficiente para sugerir alterações. Nenhuma das respostas abertas aponta erro de tradução ou de execução.



\section{Comparação A/B com estudantes}
\label{sec:avaliacao-estudantes}

Esta avaliação foi feita com uma turma de estudantes, que comparou duas versões do mesmo programa: a Versão A, na sintaxe padrão, e a Versão B, personalizada em português. O formulário também traz afirmações sobre a plataforma e a demonstração.

\subsection{Formulário e procedimento}
\label{subsec:formulario-estudantes}

O formulário dos estudantes não constava da versão anterior do protocolo. Ele foi preparado para uma sessão de demonstração com uma turma do 1º ano do ensino médio técnico em Informática da UFCI-MG, Campus Leopoldina, realizada em 3 de novembro de 2026. Os autores demonstraram a plataforma em sala, os alunos também a usaram e responderam ao formulário em seguida.

O formulário é curto. Depois do termo de assentimento, pergunta a faixa de idade, o curso, o ano e a experiência anterior com programação. Seguem seis afirmações em escala de concordância de cinco pontos: se o aluno entendeu para que serve a plataforma, se um programa escrito com palavras em português é mais fácil de entender do que um escrito com palavras em inglês, se poder escolher as palavras da linguagem ajudaria quem está começando, se ver a separação do programa em \textit{tokens} ajuda a entender como o computador lê o código, se gostaria de usar a plataforma nas aulas e se a demonstração aumentou seu interesse em aprender programação. Uma questão mostra duas versões de um mesmo programa, que lê um número e escreve se ele é par ou ímpar. A versão A usa a sintaxe padrão, em inglês, e a versão B usa uma sintaxe personalizada em português. O aluno indica qual das duas acha mais fácil de entender, ou se as duas são igualmente fáceis ou igualmente difíceis. A medida obtida com essa pergunta é a preferência declarada entre as versões, isto é, o número de alunos que escolheu cada uma das quatro alternativas; ela registra o que o aluno diz achar mais fácil e não mede compreensão. O formulário termina com um campo de comentário livre.

Dois desses itens tocam o tema deste trabalho. O item sobre os \textit{tokens} leva ao público-alvo uma pergunta sobre a exibição de uma fase intermediária da tradução, que o formulário dos especialistas só alcança pela opinião de especialistas. A questão das duas versões depende de uma propriedade que o núcleo garante, a de que vocabulários diferentes podem descrever o mesmo programa. O que ela mede, porém, é qual das duas escritas o aluno declara entender com mais facilidade.

Os limites desse instrumento são maiores que os do formulário dos especialistas. A turma é uma só e foi escolhida por conveniência. A opinião foi colhida logo depois de uma demonstração feita pelos próprios autores e de um período de uso da plataforma, situação que favorece o efeito de novidade e a tendência a concordar. O formulário não mede desempenho. Ele serve como indício de aceitação e de preferência declarada.

\subsection{Resultados da sessão}
\label{sec:resultado-estudantes}

Esta subseção trata das afirmações sobre a demonstração (questões \ref{qF:demo-ini} a \ref{qF:demo-fim}) e da pergunta de comparação entre as versões (questão \ref{qF:comparacao}) do Apêndice~\ref{ap:formulario-estudantes}. O formulário descrito na Subseção~\ref{subsec:formulario-estudantes} foi aplicado em 3 de novembro de 2026 a uma turma do 1º ano do ensino médio técnico em Informática da UFCI-MG, Campus Leopoldina, depois de uma demonstração da plataforma feita pelos autores e de um período em que os alunos também a usaram. Houve dezoito respostas, todas com o termo de assentimento marcado. Nove alunos tinham até 15 anos e nove tinham 16 ou 17. Nove disseram já ter programado um pouco, oito disseram programar com frequência e um disse nunca ter programado. Quase ninguém, portanto, era iniciante absoluto.

A Tabela~\ref{tab:resultado-estudantes} traz as seis afirmações. Todas tiveram média entre 4,00 e 4,33. A afirmação de que um programa em português é mais fácil de entender do que um em inglês teve a menor média e o maior número de neutros: seis dos dezoito não tomaram posição sobre o item mais próximo da premissa da plataforma.

\begin{table}[htbp]
\centering
\caption{Afirmações do formulário da sessão com estudantes (escala de 1 a 5; dezoito respostas).}
\label{tab:resultado-estudantes}
\footnotesize
\renewcommand{\arraystretch}{1.2}
\begin{tabular}{|p{6.3cm}|c|c|c|c|c|}
\hline
\textbf{Afirmação} & \textbf{Média} & \textbf{Mediana} & \textbf{Concordam} & \textbf{Neutros} & \textbf{Discordam} \\\hline
Entendi para que serve a plataforma & 4,33 & 4,5 & 17 & 0 & 1 \\\hline
Um programa escrito com palavras em português é mais fácil de entender do que um escrito com palavras em inglês & 4,00 & 4 & 11 & 6 & 1 \\\hline
Poder escolher as palavras da linguagem ajudaria quem está começando a programar & 4,17 & 4 & 14 & 3 & 1 \\\hline
Ver como a plataforma separa o programa em partes (os \textit{tokens}) ajuda a entender como o computador lê o código & 4,17 & 4 & 15 & 1 & 2 \\\hline
Eu gostaria de usar essa plataforma nas aulas de programação & 4,17 & 4 & 15 & 3 & 0 \\\hline
A demonstração aumentou meu interesse em aprender programação & 4,17 & 4 & 16 & 0 & 2 \\\hline
\end{tabular}
\legend{Fonte: dados da pesquisa (2026).}
\end{table}

Para este trabalho, o item central é o dos \textit{tokens}, porque leva ao público-alvo uma pergunta sobre a exibição de uma fase da tradução. Quinze dos dezoito alunos concordaram que ver a separação do programa em \textit{tokens} ajuda a entender como o computador lê o código, um ficou no ponto neutro e dois discordaram. Os oito que programam com frequência concordaram todos. Entre os nove que já tinham programado um pouco, sete concordaram, um ficou no ponto neutro e um discordou. O aluno que nunca tinha programado também discordou. No painel, a mesma fase teve média 4,50 como recurso e 4,29 como valor didático, mas só dois especialistas declararam ter visto a tabela. Os estudantes a viram na demonstração, e a concordância deles é o dado mais direto que o trabalho tem sobre essa fase. Ainda assim, o item registra concordância com uma frase e não mede compreensão.

O formulário também mostrou duas versões de um programa que lê um número e escreve se ele é par ou ímpar: a Versão A, na sintaxe padrão em inglês, e a Versão B, personalizada em português. A Tabela~\ref{tab:resultado-versoes} separa as respostas pela experiência declarada.

\begin{table}[htbp]
\centering
\caption{Versão do programa considerada mais fácil de entender, por experiência declarada com programação.}
\label{tab:resultado-versoes}
\footnotesize
\renewcommand{\arraystretch}{1.2}
\begin{tabular}{|p{4.3cm}|c|c|c|c|c|}
\hline
\textbf{Experiência} & \textbf{Versão B} & \textbf{Igualmente fáceis} & \textbf{Versão A} & \textbf{Igualmente difíceis} & \textbf{Total} \\\hline
Programa com frequência & 6 & 2 & 0 & 0 & 8 \\\hline
Já programou um pouco & 5 & 2 & 2 & 0 & 9 \\\hline
Nunca programou & 0 & 0 & 0 & 1 & 1 \\\hline
\textbf{Total} & 11 & 4 & 2 & 1 & 18 \\\hline
\end{tabular}
\legend{Fonte: dados da pesquisa (2026).}
\end{table}

Onze alunos acharam a Versão B mais fácil de entender, quatro acharam as duas igualmente fáceis, dois preferiram a Versão A e um achou as duas igualmente difíceis. A questão depende de uma propriedade que o núcleo garante: as duas versões são o mesmo programa. As palavras da Versão B são convertidas nos mesmos \textit{tokens} que as da Versão A, e a partir daí a análise sintática, o código intermediário e a execução não distinguem as duas. O formulário afirmava ao aluno que as versões fazem a mesma coisa, e o que ele mediu foi qual das duas escritas o aluno declara entender com mais facilidade. Não mediu se o aluno reconheceria a equivalência sozinho.

A preferência pela versão em português apareceu nos dois grupos com experiência. Os dois alunos que preferiram a versão em inglês estão entre os que já tinham programado um pouco, e um deles é o que discordou do item dos \textit{tokens}. O único aluno que nunca tinha programado, justamente o perfil para o qual a plataforma foi pensada, achou as duas versões igualmente difíceis e discordou de três das seis afirmações, entre elas a de ter entendido para que serve a plataforma e a dos \textit{tokens}. Uma única resposta não sustenta conclusão, mas ela vai em sentido contrário ao das demais e fica registrada.

O campo de comentário livre recebeu duas respostas: uma elogiou a plataforma e a outra pediu o acréscimo de outra linguagem. O formulário não tem item de autorização de citação, e por isso os comentários não são transcritos.

Os limites dessa sessão estão na Subseção~\ref{subsec:formulario-estudantes}. Foi uma única turma, de conveniência, que respondeu logo depois de uma demonstração feita pelos próprios autores e de um período de uso da plataforma, e quase todos já programavam. Nenhuma medida de desempenho foi colhida. As respostas servem como indício de aceitação e de preferência declarada pela versão em português. Não mostram redução da barreira sintática nem aprendizagem. O registro do assentimento e a falta de autorização dos responsáveis estão descritos na Seção~\ref{sec:protocolo-avaliacao} e retomados nas limitações do Capítulo~\ref{cap:consideracoes}.



\section{Relação com trabalhos anteriores}
\label{sec:resultado-literatura}

Esta seção compara os resultados com as ferramentas e os estudos do Capítulo~\ref{cap:referencial}. A comparação tem um limite de origem: aqui há opinião e preferência declarada, enquanto parte dos trabalhos citados mediu desempenho.

A preferência dos estudantes pela versão em português e a concordância do painel com a afirmação de que palavras em português facilitam a compreensão, com média 4,38, vão na direção da escolha feita pelo Portugol Studio, que traduz as palavras-chave para que o iniciante não precise lidar com dois idiomas ao mesmo tempo \cite{univali2026}. Os dados não alcançam uma diferença de desenho entre as duas propostas. No Portugol o vocabulário em português é fixo; na plataforma ele é escolhido por quem usa. O formulário dos estudantes comparou uma versão em inglês com uma em português definida pelos autores e, por isso, não informa o efeito de o próprio aluno escolher as palavras.

\citeonline{stefik2013} compararam o desempenho de novatos em linguagens consolidadas com o obtido em uma linguagem de controle de palavras-chave aleatórias e encontraram efeito da forma sintática sobre o desempenho. Os dados desta avaliação são de outro tipo, porque não houve tarefa nem contagem de erros. Eles não confirmam nem contrariam aquele resultado, e a comparação depende do estudo de desempenho previsto no Capítulo~\ref{cap:consideracoes}.

Cinco dos oito especialistas conheciam o Portugol Studio ou o Portugol Webstudio, e mesmo assim a afirmação de que a plataforma oferece algo que as ferramentas conhecidas não oferecem teve média 4,62. O formulário não pergunta qual seria esse diferencial. A execução passo a passo, item de maior valor didático no painel, com média 4,75, também existe no Portugol Studio, que tem depurador e inspetor de variáveis \cite{univali2026}. O que este núcleo acrescenta é a exibição dos \textit{tokens} e do código intermediário, e essas fases foram pouco visitadas: dois especialistas declararam ter visto a tabela de \textit{tokens} e quatro, o código intermediário. A avaliação diz pouco, portanto, sobre o recurso que mais distingue o núcleo das ferramentas existentes.

A correção automática pelos casos de teste, com média 4,50 entre os que lecionam, segue o ciclo de submissão e veredito de juízes \textit{online} como o Beecrowd \cite{beecrowd2026faq}. A diferença é que os casos são executados com a configuração de linguagem do aluno, como descrito na Seção~\ref{sec:integracao-plataforma}. Por fim, a preocupação com a passagem para ferramentas de mercado, que dividiu o painel no item de migração, também orienta o Portugol Studio, cuja tela reproduz a organização de uma IDE profissional para que a familiaridade seja transferível \cite{univali2026}. Lá a preocupação recai sobre a interface; aqui o painel a levantou sobre a sintaxe.

\section{Limites de interpretação e continuidade}
\label{subsec:limites-interpretacao}

Além do limite já citado sobre o público, duas ressalvas afetam a leitura dos resultados. A primeira é sobre a redação em português dos itens do SUS, para a qual não há registro de comparação com uma tradução validada. A segunda é sobre a análise: escores altos dados por especialistas não comprovam aprendizagem nem redução de carga cognitiva.

Investigar a hipótese pedagógica é a etapa seguinte e exige outro desenho de estudo: comparar, na mesma plataforma, o vocabulário padrão e o vocabulário personalizado em tarefas equivalentes de entrada e saída, seleção e repetição. A ordem das condições deve ser alternada entre os participantes, e os enunciados, os casos de teste e os critérios de classificação dos erros precisam ser definidos antes da coleta, separando erros de escrita de falhas de lógica. Com isso seria possível descrever os acertos, o tempo até a primeira execução correta e a frequência de erros em cada condição. Ainda assim, uma diferença nessas medidas não comprova, sozinha, que houve aprendizagem duradoura. Esse estudo deve ter aprovação do comitê de ética da instituição antes de qualquer coleta.
\section{Síntese e confronto com os objetivos}
\label{sec:resultado-sintese}

Esta seção retoma os objetivos específicos da introdução e diz, para cada um, o que a avaliação sustenta e o que continua sem verificação.

O primeiro objetivo era especificar a gramática da linguagem e os elementos passíveis de redefinição. O painel considerou o conjunto de elementos configuráveis adequado ao propósito da ferramenta, com média 4,57 e nenhuma nota abaixo de 4. Dois dos três professores que responderam consideraram o subconjunto suficiente para uma disciplina introdutória, e o terceiro ficou no ponto neutro. As construções apontadas como ausentes não convergiram: registros tiveram duas marcações, e as demais, uma cada.

O segundo e o terceiro objetivos eram projetar e implementar as análises léxica, sintática e semântica com suporte à configuração do vocabulário. O painel avaliou bem o resultado visível dessas fases: a resposta rápida e estável teve média 4,88, a validação prévia da configuração 4,67 e a clareza das mensagens de erro 4,71. Um respondente, que deu nota máxima às mensagens, pediu que elas sugerissem a correção, recurso que o núcleo não tem. A verificação do autor mostrou que a mensagem de \textit{token} inesperado nomeia a classe do \textit{token} pela palavra canônica, mesmo quando o usuário escreveu outra. A adaptação das mensagens ao vocabulário em uso é parcial, e esse ponto não foi captado pelo instrumento.

O quarto objetivo era integrar o núcleo à interface. A afirmação de que o código intermediário exibido corresponde ao que o programa executa teve média 4,83 entre seis respondentes, sem nenhuma nota abaixo de 4, e dois marcaram ``não sei''. Só quatro dos seis que deram nota declararam ter visto o código intermediário. O resultado não contradiz a premissa de que a interface mostra o que a máquina executa, mas é uma percepção com base pequena. A sustentação dessa premissa continua sendo a construção do sistema e os testes da Seção~\ref{sec:metodologia-testes}.

O quinto objetivo era disponibilizar o núcleo em ambiente \textit{web} e verificá-lo por testes automatizados. Cinco respondentes declararam ter escrito e executado programas na plataforma \textit{web}, e nenhum relatou erro de tradução ou de execução nas respostas abertas. A verificação por testes está descrita na Seção~\ref{sec:metodologia-testes} e não depende da avaliação.

O sexto objetivo era definir o protocolo de avaliação. O protocolo foi definido e aplicado, com oito respondentes, e a ele se somou a sessão com dezoito estudantes. A aplicação mostrou limites do próprio instrumento: dois itens mudaram de bloco e passaram a ser respondidos por no máximo quatro pessoas, o termo de consentimento ficou sem item de marcação e vários recursos receberam mais notas do que usos declarados.

Quanto à hipótese pedagógica, vale a distinção da Seção~\ref{subsec:limites-interpretacao}. O que foi coletado é opinião. A maioria do painel considera plausível que a personalização reduza a barreira sintática, com média 3,88, e um respondente discorda das três premissas. Onze de dezoito estudantes acharam a versão em português mais fácil de entender, depois de uma demonstração e de um período de uso, sem tarefa medida, e o único que nunca tinha programado achou as duas igualmente difíceis. Nenhum desses dados prova a redução da barreira nem aprendizagem. Do lado desfavorável, o painel se dividiu sobre o risco de migração para uma linguagem de mercado, três de sete respondentes concordam que a personalização pode desviar a atenção do raciocínio algorítmico e a intenção de adotar a plataforma em uma disciplina introdutória ficou em 3,00 entre os que lecionam. Esses pontos têm o mesmo peso dos resultados favoráveis e orientam os encaminhamentos do Capítulo~\ref{cap:consideracoes}.

